\runninghead{Scope and References}
\section{Appendix: Scope and Qualifications}
\label{sec:scope}

{\small
\textbf{Edition, Mass and occurrence.} The study treats the Mass
\latin{Dominica sexta quae superfuit post Epiphaniam}, II classis, of the
\work{Missale Romanum}, editio typica, Typis Polyglottis Vaticanis 1962,
printed pp.~421--422, nos.~1759--1768, as said on Sunday 15~November 2026, the
twenty-fifth Sunday after Pentecost, in the universal 1962 calendar and no
particular calendar. Its orations and lessons are those of the Sixth Sunday
after the Epiphany (pp.~48--49, nos.~337--345) and its chants those of the
Twenty-third Sunday after Pentecost, under RGMR~298. Its place in 2026, the
omission of the calendar feast of St Albert the Great, the colour, the
Gloria, the Credo and the Trinity Preface rest on the Missal's general
rubrics (RG~10, 12, 14, 18~b, 91, 107--111, 127~b; RGMR~298, 431, 434, 475,
494~b) and on the Missal's printed direction after the Secret, as
\texttt{research/context.md} records; two dated Ordos of institutes using the
1962 books, of the Fraternity of St Peter and of the Institute of Christ the
King Sovereign Priest, both in France, give the same Mass for the date. No
local oration, external solemnity or votive Mass is considered. The
\latin{Asperges}, the Ordinary and the Office are not treated.

\textbf{Text control.} Every Latin form was read on page images of the CMAA
facsimile of the 1962 typical edition and collated word for word with the
Ratisbon printing of 1862 (Pustet), whose page images carry every word of the
ten texts; the collation, the comparison with the Epiphany-season printing and
the comparison with the Clementine Vulgate are in \texttt{propers/verified.md}.
Where the chants and lessons differ from the Clementine Vulgate the
differences are reported as textual facts; the Latin psalter, antiphonary or
lectionary from which those forms came was not identified. The resumed
printing's citation of the Epistle as \latin{1 Thess.\ 1, 1-20} is a misprint
for 1~Thess 1:2--10. The rubrics of the 1962 book are described and not
reproduced.

\textbf{English.} The English of the lessons, and of each chant's canonical
verses, is the Douay--Rheims in Challoner's revision. The English of the five
chants as the Missal sings them and of the three orations is that of the
Lasance \work{New Roman Missal} in the Benziger revision of 1945, public
domain in the United States by non-renewal, which prints this Mass as its
\work{Supplementary Mass No.~1}. None of this English is an approved liturgical
translation of the 1962 Missal. Augustine's psalm expositions and Chrysostom's
homilies on Matthew and First Thessalonians are quoted in the English of the
Nicene and Post-Nicene Fathers; the expositions of Psalms~43, 77 and 84 in
that series are abridged, so that the quotations are of the condensed English
and not of Augustine's continuous text. Elsewhere the Latin of the Fathers is
quoted with an English rendering made for this study, and no English is
composed for a liturgical text.

\textbf{Reception.} The interpretations rest on the loci named in the text and
recorded, with editions and bounds, in \texttt{research/scope.md} and
\texttt{research/interpretations.md}. Every connection drawn between two texts
of the Mass belongs to the study's interpretations unless an author is named as
making it; no Father cited here comments on this Mass as a whole. Ambrose,
Chrysologus and Bede are cited on Luke's and Mark's forms of the two parables,
and Chrysostom on Matthew's form of the Communion's saying (Mt 21:21--22).
The Latin of Jerome and of Hilary on Matthew was read on page images; Ambrose
and Bede in transcriptions of Migne not collated with page images; Hilary on
the Psalms, Chrysologus, Aquinas, Rupert, Durandus and Sicard in uncorrected
optical text or locally parsed web text; Chrysostom on the Psalms and
Theodoret in uncorrected Greek optical text, so that they are given in
paraphrase. Rupert, Durandus and Sicard expound these chants in Masses whose
lessons differed from this Mass's, and they are cited for the chants alone.
Every principal author of an interpretation has a row in the repository's
register of authorial standing; Theodoret, Rupert, Durandus and Sicard are
ecclesiastical writers and are cited in support only. No commentary on the
Collect, Secret or Postcommunion survives among the commentators read.
Origen on Matthew, Cyril of Alexandria on Luke and the psalm commentaries of
Cassiodorus and Bellarmine on these chants were not read.

\textbf{Chronology.} Every biblical date in the Scriptural Date and Location
appendix is generated from the repository's Scripture chronology corpus,
with its sources named there.

\textbf{Rights.} The Latin is published on 17 U.S.C. 103(b) together with its
public-domain witness, the Pustet printing of 1862. The Douay--Rheims and the
Lasance English are quoted from registered editions. The other quotations are
short and come from editions published before 1931 or from medieval and
ancient Latin: Migne's \work{Patrologia Latina}, vols.~9, 24, 26, 52, 92, 170
and 213, and \work{Patrologia Graeca}, vols.~55 and 82; the Venice Aquinas
(1745); the Lyon Durandus (1612); and the Nicene and Post-Nicene Fathers. The
Corpus Thomisticum and Wikisource texts are used only for Latin that is itself
medieval or ancient. Scripture, liturgical texts and quoted translations are
not relicensed by the repository's licence.

\textbf{Use.} This is a study aid to the 1962 Missal, not an official
liturgical book, translation or critical edition.
\par}

\section{References}
\label{sec:references}

{\small
\subsection*{Liturgical books}
\begin{itemize}[leftmargin=1.2em,itemsep=.25em]
\item \work{Missale Romanum ex decreto Sacrosancti Concilii Tridentini
restitutum, Summorum Pontificum cura recognitum}, editio typica (Typis
Polyglottis Vaticanis, 1962), CMAA facsimile: \work{Rubricae generales} and
\work{Rubricae generales Missalis romani}; Proprium de Tempore, pp.~48--49,
nos.~337--345, and pp.~421--422, nos.~1759--1768.
\item \work{Missale Romanum} (Ratisbon: Pustet, 1862), pp.~44--45 and
357--359.
\item Lasance, F.~X., and Walsh, F.~A., \work{The New Roman Missal in Latin and
English}, revised edition (New York: Benziger Brothers, 1945), pp.~740--743.
\item Fraternité Saint-Pierre (France), \work{Ordo du mois}, and Institut du
Christ Roi Souverain Prêtre (France), \work{Ordo}, entries for 15~November
2026 (websites), read 11~November 2026.
\end{itemize}

\subsection*{Scripture}
\begin{itemize}[leftmargin=1.2em,itemsep=.25em]
\item \work{The Holy Bible Translated from the Latin Vulgate} (Douay--Rheims,
Challoner revision), Project Gutenberg text: Jer 29:1--14; Ps 43:1--26;
77:1--2; 84:1--14; 129:1--8; Mt 13:1--36; 16:6; 17:19; Mk 11:12--25;
Acts 17:1--10; 1~Cor 2:4; 5:7; 1~Thess 1:2--10; 2:2; 3:1--7.
\item \work{Biblia Sacra Vulgatae Editionis} (Clementine Vulgate), for the
comparisons with the Missal's forms.
\end{itemize}

\subsection*{Fathers, Doctors and other authors}
\begin{itemize}[leftmargin=1.2em,itemsep=.25em]
\item Ambrose, \work{Expositio evangelii secundum Lucam} VII, on Lk
13:18--21, nn.~176--194, Wikisource transcription of PL~15.
\item Augustine, \work{Enarrationes in Psalmos} 43, 77 and 84, NPNF, first
series, vol.~8; \work{Quaestiones Evangeliorum} I.11--12, Migne Latin text
(Wikisource).
\item Bede, \work{In Lucae evangelium expositio} IV, on Lk 13:18--21, PL~92,
Wikisource transcription; \work{In Marci evangelium expositio} I, on
4:30--32, and III, on 11:22--25, Migne Latin text (Wikisource).
\item Chrysologus, Peter, \work{Sermones} 98 and 99, PL~52 (Paris, 1894).
\item Chrysostom, John, \work{Homilies on Matthew} 46, 47 and 67, NPNF, first
series, vol.~10; \work{Homilies on First Thessalonians} 1 and 2, NPNF, first
series, vol.~13; \work{Expositio in Psalmos}, on Pss 43 and 129, PG~55.
\item Hilary of Poitiers, \work{In Evangelium Matthaei commentarius} 13.4--6,
PL~9 (Paris, 1844), cols.~994--995; \work{Tractatus super Psalmos}, on Ps 129,
PL~9, col.~718.
\item Jerome, \work{Commentarii in Matthaeum} II, on 13:31--35, PL~26 (Paris,
1845), cols.~90--93; \work{In Hieremiam} V, on 29:10--14, PL~24 (Paris,
1865), cols.~893--894.
\item Theodoret of Cyrus, \work{Interpretatio epistolae I ad Thessalonicenses},
on 1:3, PG~82.
\item Thomas Aquinas, \work{Super Evangelium S.~Matthaei lectura}, cap.~13, in
\work{Opera omnia}, vol.~3 (Venice: Bettinelli, 1745), pp.~186--189;
\work{Super I epistolam ad Thessalonicenses}, cap.~1, lect.~1, Corpus
Thomisticum text; \work{Super Psalmos}, on Ps 43, nn.~3--4, Corpus Thomisticum
text.
\end{itemize}

\subsection*{Commentators on the Mass}
\begin{itemize}[leftmargin=1.2em,itemsep=.25em]
\item Durandus, William, \work{Rationale divinorum officiorum} VI.140, in the
Lyon printing of 1612, tomus secundus, p.~412.
\item Rupert of Deutz, \work{De divinis officiis} XII.23, PL~170,
cols.~329--330.
\item Sicard of Cremona, \work{Mitrale} VIII.23, PL~213.
\end{itemize}
\par}
